% ============================================================
% Come si usano, argomento 3: la probabilità e le sue regole
% ============================================================
\argomento{Probabilità: definizioni e regole}{Formule a p.~\pageref*{f:frequenza}}

\begin{scelta}{Come scegliere la regola: dal testo alla formula}
\renewcommand{\arraystretch}{1.25}
\begin{tabular}{@{}p{0.47\linewidth}p{0.47\linewidth}@{}}
  \toprule
  \textbf{Se il testo\ldots} & \textbf{si usa\ldots} \\
  \midrule
  descrive risultati tutti alla pari (dado onesto, mazzo ben mescolato, scelta «a caso»)
    & casi favorevoli su casi possibili \vv{laplace}, contando con l'argomento 2 \\
  dà risultati con pesi diversi (dado truccato)
    & la somma dei pesi \vv{pesi} \\
  dice «almeno uno», «almeno un\ldots»
    & il complementare \vv{compl}: si calcola «nessuno» \\
  dice «$A$ oppure $B$»
    & l'unione \vv{unione}, togliendo l'intersezione \\
  dice «$A$ ma non $B$», «solo $A$»
    & la differenza \vv{sottrazione} \\
  dice «sapendo che\ldots», «dato che\ldots»
    & la probabilità condizionata (argomento 4) \\
  \bottomrule
\end{tabular}

Alla fine di ogni conto, un controllo: il risultato deve stare fra 0 e 1, e un evento più piccolo non può essere più
probabile di uno che lo contiene (\vv{monotonia}).
\end{scelta}

\scheda{frequenza}{
  \usoquando{Si hanno dati sperimentali o simulati: per $n$ grande la frequenza è una \emph{stima} della
    probabilità, $P(A) \simeq n_A / n$.}
  \usoesempio{Su $1\,000$ bit trasmessi ne arrivano sbagliati 12: $f_{1000}(\text{errore}) = 12/1000 = 0{,}012$,
    e la probabilità di errore del canale è circa $0{,}012$.}
}

\scheda{frequentista}{
  \usoquando{Per capire che cosa \emph{significa} una probabilità e da dove vengono le regole: ogni regola vale per
    le frequenze, e quindi per il loro limite. Non si usa per fare i conti, perché richiede infinite prove.}
  \usoattenzione{La frequenza non è una successione fissata: ogni serie di prove dà numeri diversi. Il limite va
    inteso in senso «fisico», e lo rende preciso la legge dei grandi numeri.}
}

\scheda{laplace}{
  \usoquando{Tutti i risultati dell'esperimento sono ugualmente probabili: dado o moneta onesti, carte ben mescolate,
    estrazioni «a caso».}
  \usocome{(1) Si sceglie $\Omega$ in modo che i suoi risultati siano \emph{davvero} equiprobabili (coppie ordinate
    per più lanci, sottoinsiemi per una mano di carte). (2) Si conta $\card{\Omega}$. (3) Si conta $\card{A}$ con lo
    stesso criterio (se $\Omega$ ignora l'ordine, lo ignora anche $A$). (4) Si divide.}
  \usoattenzione{L'errore di d'Alembert: con due monete i risultati «due teste», «due croci», «una testa e una croce»
    non sono equiprobabili. Quelli equiprobabili sono $TT$, $TC$, $CT$, $CC$, e $P(\text{una testa e una croce}) =
    2/4 = 1/2$, non $1/3$.}
  \usoesempio{Due lanci di un dado, «almeno un 6»: 11 coppie su 36, $P = 11/36 \simeq 0{,}306$. Poker con 32 carte,
    «esattamente due assi»: $\binom{4}{2}\binom{28}{3} / \binom{32}{5} = 19\,656 / 201\,376 \simeq 0{,}0976$.}
}

\scheda{assiomi}{
  \usoquando{Per controllare che una funzione sia davvero una probabilità, e per dimostrare le altre regole: tutte
    scrivono un evento come unione di pezzi incompatibili e applicano l'additività.}
  \usovarianti{La numerabile additività serve quando l'evento è l'unione di \emph{infiniti} eventi incompatibili.}
  \usoesempio{Si lancia una moneta onesta finché esce testa. «Prima testa al lancio $n$» ha probabilità
    $(1/2)^n$, e questi eventi sono incompatibili: $P(\text{prima o poi esce testa}) =
    \sum_{n=1}^{\infty} (1/2)^n = 1$ (\vv{geometrica}).}
}

\scheda{pesi}{
  \usoquando{I risultati non sono equiprobabili, ma si conosce la probabilità (il peso) di ciascuno.}
  \usocome{Si sommano i pesi dei risultati che fanno verificare l'evento. Prima si controlla che i pesi siano
    $\geq 0$ e sommino a 1.}
  \usoesempio{Dado truccato con $P(1) = 5/6$ e $P(i) = 1/30$ per le altre facce:
    $P(\text{dispari}) = \frac{5}{6} + \frac{1}{30} + \frac{1}{30} = \frac{9}{10}$, mentre
    $P(\text{pari}) = \frac{3}{30} = \frac{1}{10}$, anche se i due eventi hanno 3 elementi ciascuno.}
}

\scheda{compl}{
  \usoquando{Il testo dice «almeno uno», «almeno un\ldots», «non tutti»: il complementare («nessuno», «tutti») è
    quasi sempre un conto solo.}
  \usocome{Si scrive il complementare dell'evento, se ne calcola la probabilità e la si toglie da 1.}
  \usoattenzione{Il complementare di «almeno uno» è «nessuno», non «al più uno».}
  \usoesempio{Dieci lanci di un dado, «almeno un 6»: «nessun 6» ha probabilità $(5/6)^{10}$, quindi
    $P = 1 - (5/6)^{10} \simeq 0{,}8385$.}
}

\scheda{unione}{
  \usoquando{Il testo dice «$A$ oppure $B$», «almeno uno fra $A$ e $B$», e si sa calcolare l'intersezione.}
  \usocome{Si sommano le due probabilità e si toglie quella dell'intersezione, che era stata contata due volte.}
  \usovarianti{Se $A$ e $B$ sono incompatibili l'intersezione è vuota e basta sommare. Con tre eventi si somma, si
    tolgono le tre intersezioni a due a due e si rimette quella dei tre.}
  \usoattenzione{Sommare senza togliere l'intersezione può dare risultati maggiori di 1.}
  \usoesempio{Un dado, «pari oppure almeno 5»: $\frac{3}{6} + \frac{2}{6} - \frac{1}{6} = \frac{2}{3}$ (il 6 è in
    entrambi). Con tre eventi, $A = \{2, 4, 6\}$, $B = \{1, 2, 3\}$, $C = \{3, 6\}$:
    $\frac{3 + 3 + 2}{6} - \frac{1 + 1 + 1}{6} + 0 = \frac{5}{6}$, e infatti $A \cup B \cup C = \{1, 2, 3, 4, 6\}$.}
}

\scheda{sottrazione}{
  \usoquando{Il testo dice «$A$ ma non $B$», «solo $A$».}
  \usovarianti{Se $B$ è contenuto in $A$, l'intersezione è $B$ stesso e si toglie direttamente $P(B)$.}
  \usoesempio{Un dado, «pari ma non 6»: $\{6\} \subseteq \{2, 4, 6\}$, quindi
    $P = \frac{1}{2} - \frac{1}{6} = \frac{1}{3}$.}
}

\scheda{monotonia}{
  \usoquando{Come controllo di plausibilità dei risultati: ogni probabilità sta fra 0 e 1, e se $A \subseteq B$
    deve venire $P(A) \leq P(B)$.}
  \usoesempio{«Esce 6» implica «esce un pari»: $P(\{6\}) = \frac{1}{6} \leq P(\{2, 4, 6\}) = \frac{1}{2}$.}
}
